\documentclass{gescons}

\genre {Parceria}
\author{Amanda Vieira}
%\roles{Voluntária da Editares}
\title{Grupos de trabalho com a UNICIN traçam planos para fortalecer a atuação da Editares}

\begin{document}
    \makeentrevistatitle
    %\maketitle

    %\fullwidthimage{fields}{b}

    %\coverart{back/editorial}
    \coverart{../fundo-generico.png}

\begin{center}
    \hspace{0.5cm}
    \includegraphics[height=4cm]{images/Logo-Editares-com-Marca-Registrada.png}
    \hspace{1cm}
    \includegraphics[height=3cm]{images/Unicin-Logo-1024x315.png} 
\end{center}

    
    \begin{multicols}{2}



Em 2025, a Editares deu importante passo estratégico para potencializar sua atuação na produção e distribuição de gescons. Em parceria com a UNICIN e a UNIESCON, foram criados \textbf{dois grupos de trabalho (GTs)} voltados à reformulação de processos e à implantação de novas iniciativas, com foco em dois pontos principais: o fluxo editorial e a venda de livros.

Os encontros reuniram voluntários de diferentes ICs e áreas de atuação, evidenciando a importância da intercooperação na Comunidade Conscienciológica. Dessa troca nasceram dois GTs específicos: GT Editorial e GT Comercial.

\subsection*{Diagnóstico, planejamento e ação}

Durante o primeiro semestre, os subgrupos trabalharam em diagnósticos detalhados, mapeando gargalos, levantando necessidades e propondo soluções. A partir desse estudo, foram elaborados planos de ação que começaram a ser implementados no segundo semestre de 2025.

No \textbf{GT Comercial,} as iniciativas se voltaram à expansão da presença da Editares, com novos processos de marketing, estratégias de divulgação e ações de lançamento de livros. 

No \textbf{GT Editorial,} o destaque foi a criação da \emph{Escola de Editores,} um projeto voltado à capacitação de voluntários para ampliar a qualidade técnica e a celeridade na editoração de obras conscienciológicas.

\subsection*{A força do voluntariado e a união das ICs}

O movimento reforça a relevância da união entre as Instituições Conscienciocêntricas como fator-chave para a expansão da Conscienciologia. A intercooperação entre Editares, UNICIN e UNIESCON potencializa a produção e a difusão das gescons, beneficiando pesquisadores, autores e leitores.

Nesse contexto, o voluntariado se destaca como força motriz da \emph{reurbanização extrafísica} (reurbex), atuando na materialização de projetos que impactam diretamente a evolução grupal. O livro, enquanto ferramenta de autopesquisa e recomposição grupocármica, cumpre papel central no esclarecimento e na ampliação da lucidez planetária.

Com essas iniciativas, a Editares reafirma seu compromisso com a qualidade editorial, a celeridade dos processos e a ampliação da distribuição das obras conscienciológicas. Mais do que uma reorganização interna, trata-se de um movimento coletivo que coloca o livro e o conhecimento técnico ao serviço da evolução pessoal e grupal.





    \end{multicols}
\end{document}
